\chapter{Análisis y diseño del sistema}
\label{cap:analisis-diseno}

\section{Arquitectura general}
\label{sec:arquitectura}

PCT se organiza en capas de presentación, protocolo y conectividad Wi-Fi (Figura~\ref{fig:arquitectura}). El \texttt{ProtocolOrchestrator} (\texttt{PctNodeImpl}) coordina JoinManager, LinkOrchestrator, RouteOrchestrator y TopologyManager.

% FIGURA: diagrama de componentes (exportar desde docs/protocol/03-architecture.md)
\begin{figure}[htbp]
  \centering
  \fbox{\parbox{0.85\textwidth}{\centering\vspace{2cm}
    \textit{[Figura pendiente: diagrama de componentes PCT]}\\
    \small Fuente: elaboración propia basada en \texttt{03-architecture.md}
  \vspace{2cm}}}
  \caption{Arquitectura de componentes PCT}
  \label{fig:arquitectura}
\end{figure}

\section{Modo de conexión: GO homogéneo + STA legacy}
\label{sec:modo-conexion}

\subsection{Principio rector}

Todo nodo ejecuta \texttt{createGroup()} y es GO de su propio grupo. Los enlaces del árbol se establecen \textbf{exclusivamente} mediante asociación STA legacy al SSID/PSK del GO padre. Queda prohibido \texttt{WifiP2pManager.connect()} para formar topología.

\subsection{Dos interfaces lógicas en un nodo BRIDGE}

\begin{table}[htbp]
  \centering
  \caption{Interfaces lógicas del nodo BRIDGE}
  \label{tab:interfaces-bridge}
  \begin{tabular}{@{}lll@{}}
    \toprule
    \textbf{Dirección} & \textbf{Objeto de red} & \textbf{Conoce} \\
    \midrule
    Upstream   & \texttt{Network} STA (\texttt{LegacyStaRepository}) & IP gateway del padre \\
    Downstream & GO P2P propio (\texttt{GoRepository})               & IPs DHCP de hijos \\
    \bottomrule
  \end{tabular}
\end{table}

El hijo obtiene la IP del padre como gateway de la red STA (\texttt{ConnectivityManager.getLinkProperties}). El padre obtiene la IP del hijo desde la dirección remota del socket TCP entrante en su GO (típicamente 192.168.49.x).

\subsection{Secuencia bootstrap miembro}

\begin{enumerate}
  \item Scan DNS-SD \texttt{\_pct-ctrl}; lectura TXT (\texttt{nid}, \texttt{go\_ssid}, \texttt{cp}).
  \item \texttt{requestNetwork} STA al padre.
  \item \texttt{createGroup()} --- GO BRIDGE propio.
  \item \texttt{addLocalService(\_pct-ctrl)}.
  \item TCP control al gateway STA :8765 --- HELLO, JOIN\_COMMIT.
  \item Negociación canal datos :8766 vía mensajes de control.
\end{enumerate}

\section{Capa de enlace (L2)}
\label{sec:diseno-l2}

\subsection{Dos canales TCP por vecino}

Android no garantiza \emph{fairness} en un único socket TCP mezclado. PCT exige:

\begin{itemize}
  \item \textbf{Control (:8765):} HELLO, JOIN\_COMMIT, PING/PONG, TOPO\_UPDATE, señalización DATA\_CHANNEL\_*.
  \item \textbf{Datos (:8766):} exclusivamente frames DATA (\texttt{type=user}).
\end{itemize}

\subsection{NeighborRegistry}

Cada entrada almacena: \texttt{neighbor\_nid}, \texttt{iface} (UPSTREAM/DOWNSTREAM), \texttt{local\_ip}, \texttt{ctrl\_socket}, \texttt{data\_socket}, \texttt{data\_channel\_state}, rol, hop, timestamps.

Invariantes: 0--1 upstream; 0--N downstream; identidad autoritativa = \texttt{sender\_nid} en HELLO.

\subsection{Secuencia L2}

\begin{enumerate}
  \item Hijo conecta TCP :8765 tras STA connected.
  \item HELLO bidireccional.
  \item JOIN\_COMMIT (hijo $\rightarrow$ padre).
  \item DATA\_CHANNEL\_OPEN (padre $\rightarrow$ hijo).
  \item Hijo conecta TCP :8766; DATA\_CHANNEL\_ACK.
  \item TOPO\_UPDATE inicial.
\end{enumerate}

\section{Capa de red (L3)}
\label{sec:diseno-l3}

\subsection{RoutingEntry}

Campos: \texttt{dest\_nid}, \texttt{next\_hop\_nid}, \texttt{hop\_count}, \texttt{path\_seq}, \texttt{status} (ACTIVE/BACKUP/STALE/BROKEN), \texttt{last\_seen\_ms}.

\subsection{Fusión TOPO\_UPDATE}

Al recibir entrada $(dest, hop, path\_seq, origin)$: aceptar si menor \texttt{hop\_count}; en empate, mayor \texttt{path\_seq}; rechazar si \texttt{origin == self}; decrementar TTL; marcar STALE tras timeout.

\subsection{ForwardWorker}

\begin{verbatim}
FUNCIÓN forward(msg DATA):
  SI msg.dst_nid == self: ENTREGAR a aplicación
  SI hop_limit == 0 O self in path_trace: DESCARTAR
  entry ← routing_table[msg.dst_nid]
  ENVIAR por data_socket(entry.next_hop_nid)  // solo :8766
\end{verbatim}

\section{Mensajes dirigidos y difusión}
\label{sec:mensajes-broadcast}

\subsection{Unicast}

HELLO, JOIN\_COMMIT, JOIN\_OFFER, DATA dirigido por UUID destino con reenvío hop-a-hop en canal :8766.

\subsection{Difusión acotada}

\texttt{TOPO\_UPDATE} y \texttt{NODE\_DOWN} se propagan a vecinos directos con TTL (\texttt{TOPO\_TTL\_DEFAULT = 7}), sin flooding global sin control. DNS-SD actúa como broadcast de bootstrap en L1, no como canal de datos multisalto.

\section{Desconexiones en BRIDGE}
\label{sec:diseno-desconexion-bridge}

Detección: timeout PING/PONG, \texttt{Network.onLost} STA, EOF en socket upstream.

Secuencia:
\begin{enumerate}
  \item Cerrar sockets upstream; purgar NeighborRegistry y rutas vía padre.
  \item Emitir NODE\_DOWN a hijos downstream.
  \item Transitar a ISLAND; mantener GO; anunciar \texttt{\_pct-seek}.
  \item Re-JOIN; \texttt{session\_epoch} sin cambio; \texttt{epoch} red solo si ROOT caído.
\end{enumerate}

Caída solo del canal :8766: DATA\_CHANNEL\_RESET vía control, sin reiniciar GO/STA.

\section{Unión de árboles independientes}
\label{sec:union-arboles}

Dos sub-redes con ROOT y \texttt{epoch} distintos pueden fusionarse cuando un ISLAND de árbol~B recibe JOIN\_OFFER de un BRIDGE de árbol~A. El árbol absorbente incrementa \texttt{tree\_version}; las rutas cruzadas se aprenden por TOPO\_UPDATE. Anti-ciclos: \texttt{path\_trace} + \texttt{hop\_limit}.

\section{Servicios DNS-SD}
\label{sec:dns-sd-diseno}

\begin{table}[htbp]
  \centering
  \caption{Servicios DNS-SD PCT}
  \label{tab:dns-sd}
  \begin{tabular}{@{}lll@{}}
    \toprule
    \textbf{Servicio} & \textbf{Anunciante} & \textbf{Propósito} \\
    \midrule
    \texttt{\_pct-seek.\_tcp}  & ISLAND  & Intención de unión \\
    \texttt{\_pct-ctrl.\_tcp}  & ROOT, BRIDGE, LEAF & Control operativo \\
    \bottomrule
  \end{tabular}
\end{table}

\section{Máquinas de estado}
\label{sec:maquinas-estado}

Estados radio: \texttt{R\_BOOT}, \texttt{R\_GO\_READY}, \texttt{R\_GO\_UPSTREAM}, \texttt{R\_ISLAND\_SEEK}. Roles lógicos: ISLAND, ROOT, BRIDGE, LEAF. Transición crítica: pérdida upstream $\rightarrow$ ISLAND\_SEEK $\rightarrow$ re-JOIN $\rightarrow$ GO\_UPSTREAM.

\section{Síntesis del capítulo}
\label{sec:sintesis-analisis-diseno}

El diseño de PCT separa bootstrap (L1), vecindad confiable (L2) y reenvío lógico (L3), con dos objetos de red complementarios: STA para el padre y GO P2P para hijos. Las decisiones normativas prohiben atajos (connect P2P, UDP crítico, mezcla control/datos) que comprometan robustez en Android.
